\section{Principal-part negative squares and determinant form}
\label{sec:negative-square-rewrite}

The stationary residue identity of Section~\ref{sec:zero-side} admits a second
finite-dimensional description which is useful for separating the geometry of
the bad directions from the arithmetic mechanism intended to control them.
The description is exact only after the generic perturbation used in
Section~\ref{sec:multiplicity-endgame}.  Accordingly, throughout this section
let
\[
H(z)=F(z)+\eta z+\varepsilon
\]
be a real-symmetric perturbation for which all zeros of $H'$ in
$\Omega_T$ are simple, $H$ and $H'$ have no common zero, and $H'$ has no zero
on $\partial\Omega_T$.  These are precisely the properties supplied by
Lemma~\ref{lem:generic-affine}.

The point of the present section is not to add another hypothesis to the main
theorem.  It is to identify the generic negative index with a rational
generalized-Nevanlinna defect and then with a concrete determinant zero count.
This gives a canonical scalar encoding of the quantity that the right-edge
four-block decomposition must ultimately control.

\subsection{The local principal part of $H/H'$}

Put
\begin{equation}
M(z):=\frac{H(z)}{H'(z)}.
\label{eq:full-M-H-over-Hprime}
\end{equation}
Let the real zeros of $H'$ in $\Omega_T$ be
$c_1,\dots,c_r$, and let the nonreal zeros be
\[
\rho_1,\bar\rho_1,\dots,\rho_p,\bar\rho_p,
\qquad \Im\rho_k>0.
\]
Since all stationary zeros are simple and none is a zero of $H$, define
\[
a_j:=\Res_{z=c_j}M(z)=\frac{H(c_j)}{H''(c_j)}
\in\mathbb R\setminus\{0\},
\]
and
\[
b_k:=\Res_{z=\rho_k}M(z)=\frac{H(\rho_k)}{H''(\rho_k)}
\in\mathbb C\setminus\{0\}.
\]
Real symmetry gives the residue $\overline{b_k}$ at $\bar\rho_k$.

Define the rational principal part associated with the current good rectangle by
\begin{equation}
\boxed{
M_{\Omega}(z)
:=
\sum_{j=1}^{r}\frac{a_j}{z-c_j}
+
\sum_{k=1}^{p}
\left(
\frac{b_k}{z-\rho_k}
+
\frac{\overline{b_k}}{z-\bar\rho_k}
\right).
}
\label{eq:M-Omega-principal-part}
\end{equation}

\begin{proposition}[Exact principal-part reduction]
\label{prop:principal-part-reduction}
For every pair of packet test functions $h_1,h_2$,
\begin{equation}
\boxed{
\mathcal A_H(h_1,h_2)
=
-\frac{1}{2\pi iL^2}
\oint_{\partial\Omega_T}
M_{\Omega}(z)
\,[H''h_2](z)\,[H''h_1]^\#(z)\,dz.
}
\label{eq:principal-part-contour}
\end{equation}
In particular, the holomorphic background of $H/H'$ relative to
$\Omega_T$ has exactly zero contribution to the closed contour form.
\end{proposition}

\begin{proof}
Equation~\eqref{eq:zero-kernel}, with $F$ replaced by $H$, gives
\[
\mathcal A_H(h_1,h_2)
=
-\frac{1}{2\pi iL^2}
\oint_{\partial\Omega_T}
M(z)
\,[H''h_2](z)\,[H''h_1]^\#(z)\,dz.
\]
By construction $M-M_\Omega$ is holomorphic on a neighbourhood of
$\overline{\Omega_T}$.  The packet factors and $H''$ are holomorphic there as
well.  Cauchy's theorem therefore removes the $M-M_\Omega$ term exactly.
\end{proof}

The restriction to the principal part in
\eqref{eq:M-Omega-principal-part} is essential.  A holomorphic background is
invisible to the closed contour but need not be invisible to a global Pick
kernel.  Thus no generalized-Nevanlinna assertion is made for the full
function $H/H'$.

\subsection{Negative squares}

For $z,w$ in the upper half-plane away from the poles of $M_\Omega$, define
\begin{equation}
K_{M_\Omega}(z,w)
:=
\frac{M_\Omega(z)-\overline{M_\Omega(w)}}{z-\bar w}.
\label{eq:MOmega-Nevanlinna-kernel}
\end{equation}

For a real pole one has
\begin{equation}
\frac{a_j/(z-c_j)-a_j/(\bar w-c_j)}
{z-\bar w}
=
-\frac{a_j}{(z-c_j)(\bar w-c_j)}.
\label{eq:real-pole-Pick-block}
\end{equation}
For a nonreal conjugate pair,
\begin{align}
&\frac{1}{z-\bar w}
\left[
\frac{b_k}{z-\rho_k}
+\frac{\overline{b_k}}{z-\bar\rho_k}
-\frac{\overline{b_k}}{\bar w-\bar\rho_k}
-\frac{b_k}{\bar w-\rho_k}
\right]
\notag\\
&\qquad=
-\frac{b_k}{(z-\rho_k)(\bar w-\rho_k)}
-\frac{\overline{b_k}}
{(z-\bar\rho_k)(\bar w-\bar\rho_k)}.
\label{eq:complex-pole-Pick-block}
\end{align}
In the Cauchy basis
\[
\frac1{z-\rho_k},
\qquad
\frac1{z-\bar\rho_k},
\]
the coefficient matrix of
\eqref{eq:complex-pole-Pick-block} is
\begin{equation}
\begin{pmatrix}
0&-b_k\\
-\overline{b_k}&0
\end{pmatrix},
\label{eq:complex-Pick-signature-block}
\end{equation}
whose eigenvalues are $\pm|b_k|$.

\begin{proposition}[Exact negative-square encoding]
\label{prop:negative-square-encoding}
The kernel $K_{M_\Omega}$ has exactly
\begin{equation}
\boxed{
\kappa_T(H):=B(H)+P(H')
}
\label{eq:kappa-generic-def}
\end{equation}
negative squares.  Consequently, in the scalar generalized-Nevanlinna
notation,
\[
M_\Omega\in N_{\kappa_T(H)},
\]
and
\begin{equation}
\boxed{
n_-(A_H)
=
\operatorname{sq}_-(K_{M_\Omega})
=
\kappa_T(H).
}
\label{eq:AH-negative-squares}
\end{equation}
\end{proposition}

\begin{proof}
Collect the Cauchy functions associated with the distinct poles into a vector
$\Phi(z)$.  Equations~\eqref{eq:real-pole-Pick-block} and
\eqref{eq:complex-Pick-signature-block} give a finite-rank factorization
\[
K_{M_\Omega}(z,w)
=
\Phi(z)\mathcal J_\Omega\Phi(w)^*,
\]
where
\[
\mathcal J_\Omega
=
\bigoplus_{j=1}^{r}[-a_j]
\oplus
\bigoplus_{k=1}^{p}
\begin{pmatrix}
0&-b_k\\
-\overline{b_k}&0
\end{pmatrix}.
\]
The Cauchy functions associated with distinct poles are linearly independent:
a linear relation on an open set is a rational identity, and taking residues
at the distinct poles makes every coefficient vanish.  Hence one can choose a
finite set of upper-half-plane sample points on which the Cauchy evaluation
matrix has full column rank.  Sylvester inertia then shows that the maximal
number of negative eigenvalues of a kernel Gram matrix equals the negative
inertia of $\mathcal J_\Omega$.

A real block $[-a_j]$ is negative exactly when $a_j>0$.  Since
\[
a_j=\frac{H(c_j)}{H''(c_j)},
\]
this is equivalent to $H(c_j)H''(c_j)>0$, namely to a bad real stationary
point.  Each block \eqref{eq:complex-Pick-signature-block} contributes exactly
one negative direction.  Thus the negative-square index is
$B(H)+P(H')$.  Equation~\eqref{eq:zero-inertia-generic} gives the same number
for $n_-(A_H)$.
\end{proof}

This is the precise sense in which the zero-side defect is a counting
invariant.  The magnitudes of the negative eigenvalues are irrelevant: changing
a nonzero residue without crossing the signature-changing locus does not change
the corresponding contribution to $\kappa_T(H)$.

\subsection{Cayley transform and Blaschke degree}

Define
\begin{equation}
S_\Omega(z)
:=
\frac{M_\Omega(z)-i}{M_\Omega(z)+i}.
\label{eq:SOmega-Cayley}
\end{equation}
Its upper-half-plane Schur kernel is
\[
K_{S_\Omega}(z,w)
:=
\frac{1-S_\Omega(z)\overline{S_\Omega(w)}}
{-i(z-\bar w)}.
\]
A direct calculation gives
\begin{equation}
\boxed{
K_{S_\Omega}(z,w)
=
\frac{
2K_{M_\Omega}(z,w)
}{
(M_\Omega(z)+i)
(\overline{M_\Omega(w)}-i)
}.
}
\label{eq:Cayley-kernel-congruence}
\end{equation}
Thus every finite Gram matrix for $K_{S_\Omega}$ is obtained from the
corresponding Gram matrix for $K_{M_\Omega}$ by an invertible diagonal
congruence and multiplication by the positive scalar $2$, away from the poles.
Therefore
\begin{equation}
\operatorname{sq}_-(K_{S_\Omega})
=
\kappa_T(H).
\label{eq:Cayley-same-negative-index}
\end{equation}

After conformally mapping the upper half-plane to the unit disk and, if
necessary, precomposing with a disk automorphism to choose a regular base
point, $S_\Omega$ is a scalar generalized Schur function of index
$\kappa_T(H)$.  The scalar Kre\u{\i}n--Langer factorization then gives a
coprime representation
\begin{equation}
\widehat S_\Omega
=
B_\Omega^{-1}S_0,
\label{eq:Krein-Langer-factorization}
\end{equation}
where $S_0$ is an ordinary Schur function and $B_\Omega$ is a finite Blaschke
product of degree
\begin{equation}
\boxed{
\deg B_\Omega=\kappa_T(H).
}
\label{eq:Blaschke-degree-kappa}
\end{equation}
We use only this scalar Hilbert-space case of the standard
Kre\u{\i}n--Langer theorem; see \cite{Lilleberg2020}.  In particular,
\begin{equation}
\boxed{
n_-(A_H)
=
B(H)+P(H')
=
\operatorname{sq}_-(K_{M_\Omega})
=
\operatorname{sq}_-(K_{S_\Omega})
=
\deg B_\Omega.
}
\label{eq:generic-equivalence-chain}
\end{equation}

\subsection{Polynomial and determinant form}

List all stationary poles in $\Omega_T$ as
\[
\lambda_1,\dots,\lambda_m,
\qquad
m=r+2p,
\]
and write their residues as $r_1,\dots,r_m$, so that
\begin{equation}
M_\Omega(z)
=
\sum_{j=1}^{m}\frac{r_j}{z-\lambda_j}.
\label{eq:MOmega-sum-all-poles}
\end{equation}
Put
\begin{equation}
Q_\Omega(z)
:=
\prod_{j=1}^{m}(z-\lambda_j),
\qquad
P_\Omega(z)
:=
\sum_{j=1}^{m}
r_j\prod_{\ell\ne j}(z-\lambda_\ell).
\label{eq:PQ-Omega}
\end{equation}
Then
\[
M_\Omega=\frac{P_\Omega}{Q_\Omega},
\qquad
S_\Omega
=
\frac{P_\Omega-iQ_\Omega}
{P_\Omega+iQ_\Omega}.
\]
Since every $r_j$ is nonzero, $P_\Omega$ and $Q_\Omega$ have no common zero.
Hence $P_\Omega+iQ_\Omega$ and $P_\Omega-iQ_\Omega$ have no common zero
either.  The poles of $S_\Omega$ in the upper half-plane are therefore exactly
the upper-half-plane zeros of $P_\Omega+iQ_\Omega$, counted with multiplicity.
By \eqref{eq:Blaschke-degree-kappa},
\begin{equation}
\boxed{
\kappa_T(H)
=
N_{\mathbb C_+}(P_\Omega+iQ_\Omega).
}
\label{eq:polynomial-zero-count-kappa}
\end{equation}

There is also an exact rank-one determinant representation.  Let
\[
D_\Omega:=\operatorname{diag}(\lambda_1,\dots,\lambda_m),
\qquad
r_\Omega:=(r_1,\dots,r_m)^T,
\qquad
\mathbf 1:=(1,\dots,1)^T.
\]
Then
\[
M_\Omega(z)
=
\mathbf 1^T(zI-D_\Omega)^{-1}r_\Omega.
\]
The matrix determinant lemma gives
\begin{align}
\det(zI-D_\Omega-i r_\Omega\mathbf1^T)
&=
\det(zI-D_\Omega)
\left(
1-i\,\mathbf1^T(zI-D_\Omega)^{-1}r_\Omega
\right)
\notag\\
&=
Q_\Omega(z)\bigl(1-iM_\Omega(z)\bigr).
\label{eq:determinant-lemma-Omega}
\end{align}
Since
\[
P_\Omega+iQ_\Omega
=
iQ_\Omega(1-iM_\Omega),
\]
we obtain
\begin{equation}
\boxed{
P_\Omega(z)+iQ_\Omega(z)
=
i\det(zI-D_\Omega-i r_\Omega\mathbf1^T).
}
\label{eq:determinant-polynomial-identity}
\end{equation}
Consequently,
\begin{equation}
\boxed{
\kappa_T(H)
=
\#\left\{
\lambda\in\spec(D_\Omega+i r_\Omega\mathbf1^T):
\Im\lambda>0
\right\},
}
\label{eq:rank-one-half-plane-count}
\end{equation}
with algebraic multiplicity.

The matrix in \eqref{eq:rank-one-half-plane-count} is diagonal plus rank one,
but it is non-selfadjoint.  Rank one alone therefore gives no bound on the
number of eigenvalues that can move into the upper half-plane.  The
determinant form is an exact rewrite, not a solution of the density problem.

\subsection{Interface with the arithmetic four-block decomposition}

It is important to distinguish the exact generic equivalence above from the
unperturbed arithmetic matrix.  Proposition~\ref{prop:exact-four-block-decomposition}
states exactly
\begin{equation}
A_F
=
A_T^{\rm true}
+L_T^{\rm geom}
+J_T^{\rm arch}
+R_T^{\rm exact}.
\label{eq:four-block-restated-negative-square-section}
\end{equation}
The matrix $A_F$ is not replaced by the determinant
\eqref{eq:determinant-polynomial-identity}.  Instead, Section~\ref{sec:multiplicity-endgame}
chooses $H$ on the same finite packet space so that $A_H-A_F$ is arbitrarily
small in the normalized operator metric required by the min--max argument.
Only after this genericization does
\eqref{eq:generic-equivalence-chain} apply exactly.

Thus the two sides of the method have distinct roles:
\[
\boxed{
\begin{array}{ll}
\text{zero side:}
&
A_H\longleftrightarrow
M_\Omega\longleftrightarrow
\deg B_\Omega
\longleftrightarrow
N_{\mathbb C_+}(P_\Omega+iQ_\Omega),
\\[1mm]
\text{arithmetic side:}
&
A_F=
A_T^{\rm true}
+L_T^{\rm geom}
+J_T^{\rm arch}
+R_T^{\rm exact}.
\end{array}
}
\]
The first line identifies the defect to be counted.  The second line is the
available right-edge representation from which that defect must be controlled.

In particular, the four-block decomposition is not superseded by the
determinant rewrite.  What is discarded is only the interpretation of
$A_T^{\rm true}$ as a separately positive principal matrix without controlling
its coupling to the other exact blocks.  The finite computations in
\texttt{numerical.tex} already show that the inertia can change substantially
as the missing blocks are restored.

The manuscript's existing generic count
\eqref{eq:D-H} and \eqref{eq:generic-equivalence-chain} give
\begin{equation}
\boxed{
D(H)
=
\deg B_\Omega+o(N_T).
}
\label{eq:DH-Blaschke-degree}
\end{equation}
This is an exact structural reformulation up to the already established
good-height counting error.  It is not a new estimate, because
$M_\Omega$ was built from the stationary poles themselves.

The genuinely new arithmetic problem suggested by the rewrite is therefore:
can the complete right-edge data in
\eqref{eq:four-block-restated-negative-square-section} force
\begin{equation}
\boxed{
\deg B_\Omega=o(N_T)
}
\label{eq:Blaschke-target}
\end{equation}
without first reconstructing all stationary poles?  Equivalently, can one
show from the arithmetic side that, after deleting only $o(N_T)$ exceptional
directions, the local principal part has ordinary Herglotz--Nevanlinna sign
structure?  This question is not resolved in the present paper.

The use of negative inertia to encode failures of zeta-zero positivity has
precedents in the Hermitian-form and Weil-functional literature.  Yoshida
studied Hermitian forms attached to zeta functions, while Bombieri showed for
finite truncations of Weil's quadratic functional that, under a finiteness
assumption on off-line zeros, the eventual number of negative eigenvalues is
one half of the number of zeros violating RH
\cite{Yoshida1992,Bombieri2000}.  The present principal-part construction is
different: it is tied to the stationary detector $H'$ and to the exact
Hardy-curvature contour form, and its role here is to expose the integer defect
which the four-block right-edge representation would need to control.
